\section{Worked problems: counting a step}

\subsection{Sixty-four parameters short}
\textbf{Problem.} Count Llama~3.2~3B's parameters from its shape alone
(28 layers, $d = 3{,}072$, $H_q = 24$, $H_{kv} = 8$, $d_h = 128$,
$d_{\mathrm{ff}} = 8{,}192$, $V = 128{,}256$, tied tables, one RMSNorm gain
vector before each block and one at the end) and compare with the file's
3,212,749,888. Account for any difference. Then say what untying the tables
would cost in memory and in bytes read per token.

\textbf{Solution.} Per layer, equation~\eqref{eq:params-attn} gives
$3{,}072 \times 3{,}072 = 9{,}437{,}184$ for $\mathbf{W}_Q$,
$2 \times 3{,}072 \times 1{,}024 = 6{,}291{,}456$ for $\mathbf{W}_K$ and
$\mathbf{W}_V$, and 9,437,184 for $\mathbf{W}_O$: 25,165,824. The FFN is
$3 \times 3{,}072 \times 8{,}192 = 75{,}497{,}472$, so the matrices of one layer
hold 100,663,296 and the two gains 6,144 more. Twenty-eight layers:
$2{,}818{,}572{,}288 + 172{,}032$. The tied table adds
$128{,}256 \times 3{,}072 = 394{,}002{,}432$ once, and the final norm 3,072.
Total: 3,212,749,824 --- exactly 64 short. The file's tensor list names the
missing tensor: \texttt{rope\_freqs.weight}, 64 floats, one frequency factor for
each of the $d_h/2 = 64$ rotary pairs (the chapter's section on position reads
them). No equation of the chapter predicts it, which is the point of checking a
derivation against the file rather than against another derivation.
Untying the tables would add a second $128{,}256 \times 3{,}072$ matrix:
394~million parameters, 12\% more, and 323~MB more in this file's
\texttt{Q6\_K}. It would add almost nothing per token: the output projection
reads the whole table every step either way, and the input side reads one row.
Tying saves capacity, not bandwidth.

\subsection{When the cache overtakes the weights}
\textbf{Problem.} At what context length $S$ does attention over the cache
cost as many FLOPs as the rest of a batch-one decode step? As many bytes? Answer
for Llama~3.2~3B (the \texttt{Q4\_K\_M} file, 16-bit cache) and for Qwen3-8B in
BF16 and in its \texttt{Q4\_K\_M} file, and explain why one crossover depends on
the file and the other does not.

\textbf{Solution.} Attention over the cache costs $4\,L\,H_q\,d_h$ FLOPs and
$2\,L\,H_{kv}\,d_h\,b_{kv}$ bytes per cached token; the rest of the step costs
$2P_{\text{act}}$ FLOPs and the weights' bytes. For Llama~3.2~3B:
$4 \times 28 \times 24 \times 128 = 344{,}064$ FLOPs per cached token against
$6.425$~GFLOP, so $S = 18{,}674$; and $114{,}688$ bytes per cached token against
2.01~GB of weights, so $S = 17{,}539$. For Qwen3-8B: $589{,}824$ FLOPs per
cached token against $2 \times (8.191 - 0.622)$~billion $= 15.14$~GFLOP, so
$S = 25{,}663$; and $147{,}456$ bytes per cached token against 15.14~GB of BF16
weights read, $S = 102{,}653$, or against 4.87~GB in the \texttt{Q4\_K\_M}
file, $S = 33{,}022$ (all derived). The FLOP crossover is a property of the
architecture: precision changes neither side. The byte crossover compares two
storage formats. Quantizing the weights to about 4.8 bits while the cache stays
at 16 moved Qwen3-8B's byte crossover from about 103,000 tokens to about
33,000: past that length the cache, not the weights, is what a decode step
mostly reads, and the next thing to shrink (\Chref{kv-quantization},
\Chref{shrinking-kv}). In Llama's file the two crossovers nearly coincide,
so at long context both the arithmetic and the bytes of a step are mostly
attention.

\subsection{When a mixture of experts stops saving bandwidth}
\textbf{Problem.} Granite~3.1~3B-A800M's file holds 1,828~MB of routed-expert
weights (40 experts per layer, 8 per token) and 205~MB of everything else. If
each token's experts were drawn uniformly at random, how many bytes would a
step read, per token, for batches of 1, 4 and 16 tokens? Compare with
Llama~3.2~3B's 2,012~MB, dense, and find the batch size at which Granite reads
90\% of its experts. What advantage does Granite keep once it has lost this
one?

\textbf{Solution.} A batch of $B$ tokens touches a fraction $1 - 0.8^B$ of each
layer's experts: 20\%, 59\% and 97.2\% for $B = 1$, 4 and 16. A step then reads
$205 + 1{,}828 \times (1 - 0.8^B)$~MB: 571, 1,284 and 1,982~MB, which is 571,
321 and 124~MB per token. Llama reads 2,012~MB per step at every batch size:
2,012, 503 and 126~MB per token. At batch one Granite reads 3.5 times fewer
bytes per token; at sixteen the two are level (all derived from the files).
Ninety percent of the experts are touched once $0.8^B \le 0.1$, at
$B = \ln 0.1 / \ln 0.8 = 10.3$, so eleven tokens (91.4\%); for Qwen3-30B-A3B,
with 8 of 128, it takes 36. What Granite keeps is compute: each token still
multiplies only its own eight experts, 1.77~GFLOP per token against Llama's
6.43 at short context. A mixture of experts trades bandwidth for capacity at
small batches and compute for capacity at large ones, which is why the same
model can be a laptop's fastest option at batch one and a data center's
cheapest at batch several hundred (\Chref{mixture-of-experts},
\Chref{expert-parallelism}). Routing that favours some experts makes the
union grow more slowly than this uniform case, and the advantage last longer;
\Chref{mixture-of-experts} simulates both.
